% =============================================================================
% TinyKernel — Documento Fundador
% =============================================================================
% document_id: TK-FND-00
% title: TinyKernel — Fundamentos para uma Ciência Experimental da Minimalidade Causal
% subtitle: Do “menor conjunto de componentes” ao “menor circuito causal que preserva um fenômeno”
% version: 0.1.0
% status: BASE FUNDADORA / PRÉ-IMPLEMENTAÇÃO / FALSIFICÁVEL
% project: TinyKernel
% family: TinyLogic
% author: José Pedro Trindade
% date: 2026-09-09
% language: pt-BR
% source_primary: diálogo de pesquisa “Recuperar essência do Kernel”, 2026-09-09
% source_lineage: Sister-Kernel; TinyLogicLM; TinyLogicVision
% epistemic_scope: hipóteses, definições operacionais candidatas e programa experimental
% evidence_status: nenhuma hipótese deste documento deve ser tratada como resultado experimental
% implementation_authority: NONE
% sister_kernel_authority: NONE
% motto: Sempre pronto. Sempre incompleto.
% =============================================================================

\documentclass[11pt,a4paper]{article}

\usepackage[T1]{fontenc}
\usepackage[utf8]{inputenc}
\usepackage[brazil]{babel}
\usepackage{lmodern}
\usepackage{microtype}
\usepackage{geometry}
\geometry{top=22mm,bottom=24mm,left=24mm,right=24mm}

\usepackage{amsmath,amssymb,mathtools}
\usepackage{booktabs,longtable,tabularx,array}
\usepackage{enumitem}
\usepackage{xcolor}
\usepackage{graphicx}
\usepackage{tikz}
\usetikzlibrary{arrows.meta,positioning,calc,fit,shapes.geometric}
\usepackage[most]{tcolorbox}
\usepackage{listings}
\usepackage{fancyhdr}
\usepackage{titlesec}
\usepackage{hyperref}

% -----------------------------------------------------------------------------
% Palette
% -----------------------------------------------------------------------------
\definecolor{TKInk}{HTML}{16202A}
\definecolor{TKBlue}{HTML}{2457A6}
\definecolor{TKTeal}{HTML}{157A7A}
\definecolor{TKGold}{HTML}{B27A14}
\definecolor{TKRed}{HTML}{A8423C}
\definecolor{TKGreen}{HTML}{3C7A57}
\definecolor{TKGray}{HTML}{F2F4F6}
\definecolor{TKMidGray}{HTML}{66727E}
\definecolor{TKDark}{HTML}{20262D}

\hypersetup{
  colorlinks=true,
  linkcolor=TKBlue,
  urlcolor=TKTeal,
  citecolor=TKBlue,
  pdftitle={TinyKernel — Fundamentos para uma Ciência Experimental da Minimalidade Causal},
  pdfauthor={José Pedro Trindade},
  pdfsubject={Documento fundador do laboratório TinyKernel},
  pdfkeywords={TinyKernel, minimalidade causal, kernel, causalidade, redução experimental, witness, ablação, TinyLogic}
}

% -----------------------------------------------------------------------------
% Metadata commands
% -----------------------------------------------------------------------------
\newcommand{\TKProject}{TinyKernel}
\newcommand{\TKDocID}{TK-FND-00}
\newcommand{\TKVersion}{0.1.0}
\newcommand{\TKDate}{09 de setembro de 2026}
\newcommand{\TKStatus}{BASE FUNDADORA / PRÉ-IMPLEMENTAÇÃO}
\newcommand{\TKMotto}{Sempre pronto. Sempre incompleto.}

% -----------------------------------------------------------------------------
% Layout
% -----------------------------------------------------------------------------
\pagestyle{fancy}
\fancyhf{}
\lhead{\small\sffamily\TKProject}
\rhead{\small\sffamily\TKDocID\ — v\TKVersion}
\cfoot{\small\sffamily\thepage}

\titleformat{\section}
  {\Large\bfseries\sffamily\color{TKInk}}
  {\thesection}{0.75em}{}
\titleformat{\subsection}
  {\large\bfseries\sffamily\color{TKInk}}
  {\thesubsection}{0.75em}{}
\titleformat{\subsubsection}
  {\normalsize\bfseries\sffamily\color{TKInk}}
  {\thesubsubsection}{0.75em}{}

\setlength{\parindent}{0pt}
\setlength{\parskip}{0.65em}
\setlist[itemize]{topsep=0.35em,itemsep=0.25em}
\setlist[enumerate]{topsep=0.35em,itemsep=0.25em}

% -----------------------------------------------------------------------------
% Boxes
% -----------------------------------------------------------------------------
\newtcolorbox{thesisbox}[1][]{
  enhanced,
  breakable,
  colback=TKBlue!5,
  colframe=TKBlue!70,
  boxrule=0.7pt,
  arc=2mm,
  left=3mm,right=3mm,top=2.5mm,bottom=2.5mm,
  title=#1,
  fonttitle=\bfseries\sffamily
}

\newtcolorbox{warningbox}[1][]{
  enhanced,
  breakable,
  colback=TKGold!7,
  colframe=TKGold!70,
  boxrule=0.7pt,
  arc=2mm,
  left=3mm,right=3mm,top=2.5mm,bottom=2.5mm,
  title=#1,
  fonttitle=\bfseries\sffamily
}

\newtcolorbox{falsifybox}[1][]{
  enhanced,
  breakable,
  colback=TKRed!5,
  colframe=TKRed!70,
  boxrule=0.7pt,
  arc=2mm,
  left=3mm,right=3mm,top=2.5mm,bottom=2.5mm,
  title=#1,
  fonttitle=\bfseries\sffamily
}

\newtcolorbox{markerbox}[1][]{
  enhanced,
  breakable,
  colback=TKGray,
  colframe=TKMidGray!65,
  boxrule=0.6pt,
  arc=1.5mm,
  left=3mm,right=3mm,top=2mm,bottom=2mm,
  title=#1,
  fonttitle=\bfseries\sffamily
}

% -----------------------------------------------------------------------------
% Listings
% -----------------------------------------------------------------------------
\lstdefinestyle{tk}{
  basicstyle=\ttfamily\small,
  backgroundcolor=\color{TKGray},
  frame=single,
  rulecolor=\color{TKMidGray!40},
  framerule=0.4pt,
  xleftmargin=0.5em,
  framexleftmargin=0.5em,
  breaklines=true,
  showstringspaces=false,
  columns=fullflexible
}
\lstset{style=tk}

% =============================================================================
\begin{document}

% =============================================================================
% Title page
% =============================================================================
\begin{titlepage}
\pagecolor{TKDark}
\color{white}

\vspace*{1.5cm}

{\sffamily\small
\textbf{\TKProject}\par
Família experimental TinyLogic\par}

\vspace{2.0cm}

{\sffamily\bfseries\fontsize{28}{33}\selectfont
Fundamentos para uma Ciência Experimental da Minimalidade Causal\par}

\vspace{0.6cm}

{\sffamily\fontsize{14}{19}\selectfont\color{white!78}
Do “menor conjunto de componentes” ao “menor circuito causal que preserva um fenômeno”\par}

\vspace{1.2cm}

\begin{tcolorbox}[
  enhanced,
  colback=white!6,
  colframe=white!35,
  boxrule=0.6pt,
  arc=2mm,
  left=4mm,right=4mm,top=3mm,bottom=3mm
]
\color{white}
\textbf{Proposição de partida.}
Encontrar um Kernel não consiste em escolher antecipadamente seus componentes,
mas em remover abstrações até que uma remoção adicional destrua causalidade,
significado ou capacidade.
\end{tcolorbox}

\vfill

{\sffamily
\begin{tabular}{@{}ll@{}}
Documento & \TKDocID \\
Versão & \TKVersion \\
Status & \TKStatus \\
Autor & José Pedro Trindade \\
Data & \TKDate \\
\end{tabular}
}

\vspace{1.1cm}

{\sffamily\bfseries\color{white!92}\TKMotto\par}

\end{titlepage}
\nopagecolor
\color{TKInk}

% =============================================================================
\section*{Metadados e autoridade do documento}
\addcontentsline{toc}{section}{Metadados e autoridade do documento}

\begin{tabularx}{\textwidth}{>{\bfseries}p{4.1cm}X}
\toprule
Projeto & TinyKernel \\
Documento & TK-FND-00 \\
Versão & 0.1.0 \\
Estado & Base fundadora, pré-implementação, deliberadamente falsificável \\
Fonte primária & Diálogo de pesquisa “Recuperar essência do Kernel”, consolidado em 09/09/2026 \\
Linhagem conceitual & Sister-Kernel; TinyLogicLM; TinyLogicVision \\
Objeto candidato & Minimalidade causal \\
Objeto que \emph{não} deve ser presumido & “Kernel” como entidade pronta, universal ou previamente definida \\
Autoridade de implementação & Nenhuma. Este documento não autoriza código nem experimento por si só \\
Autoridade sobre Sister-Kernel & Nenhuma. Resultados de TinyKernel só podem ser transferidos posteriormente mediante evidência e decisão própria \\
Status epistemológico & Hipóteses, definições operacionais candidatas, distinções e programa experimental; não contém resultados experimentais de TinyKernel \\
Princípio operacional & \textbf{READY $\neq$ COMPLETE} \\
\bottomrule
\end{tabularx}

\begin{warningbox}[Leitura obrigatória]
Este texto não deve ser usado como especificação de uma arquitetura que precisa ser
confirmada. Sua função é construir um aparato suficientemente pequeno para que as
próprias proposições aqui registradas possam falhar.

Se o laboratório apenas reencontrar a teoria que o criou, ele terá produzido
circularidade, não descoberta.
\end{warningbox}

\tableofcontents
\newpage

% =============================================================================
\section{Natureza deste documento}

Este é o registro fundador do \TKProject, um laboratório experimental proposto a
partir de uma convergência entre três trajetórias:

\begin{enumerate}
  \item a busca do Sister-Kernel por um conjunto mínimo de primitivas, invariantes e
  mecanismos capazes de preservar identidade e evolução;
  \item a materialização, em TinyLogicLM, de uma cadeia inteira de aprendizado em
  escala humana;
  \item a extensão, em TinyLogicVision, da pergunta sobre aprendizado para a relação
  entre mundo observado, suporte, representação, inferência, decisão e proveniência.
\end{enumerate}

A hipótese que nasce dessa convergência não é simplesmente que sistemas complexos
possuem um ``núcleo''. A hipótese mais forte é que talvez possamos investigar
experimentalmente \emph{o que precisa permanecer causalmente verdadeiro para que um
fenômeno continue sendo reconhecível como ele mesmo}.

O laboratório existe para testar essa hipótese.

\begin{thesisbox}[Tese de abertura]
\textbf{Kernel não será tratado inicialmente como objeto dado.}
O objeto científico do laboratório será a \textbf{minimalidade causal}. Um Kernel,
se essa categoria sobreviver à investigação, será um resultado candidato obtido sob
um fenômeno, contexto, família de intervenções e conjunto de testemunhos explicitados.
\end{thesisbox}

% =============================================================================
\section{A mudança de pergunta}

A formulação inicial de Kernel tendia naturalmente a perguntas do tipo:

\begin{center}
\emph{Quais são as entidades mínimas?}\\
\emph{Quais são as relações mínimas?}\\
\emph{Quais são os invariantes mínimos?}
\end{center}

Essas perguntas continuam úteis, mas carregam uma premissa oculta: a de que o Kernel
é uma coleção de coisas.

TinyLogicLM sugere uma pergunta mais fundamental:

\begin{thesisbox}[Mudança central]
Não perguntar primeiro \textbf{“qual é o menor conjunto de componentes?”}, mas:

\begin{center}
\Large\bfseries
“qual é o menor circuito causal que ainda produz o fenômeno?”
\end{center}
\end{thesisbox}

Essa mudança desloca o foco de substantivos para relações e transformações.

Um Kernel pode conter componentes, mas sua identidade pode residir menos nos
componentes isolados do que no fechamento causal entre eles.

% =============================================================================
\section{Linhagem: o que TinyLogic tornou visível}

\subsection{TinyLogicLM: capacidade como fechamento causal}

TinyLogicLM tornou observável, em uma escala em que cada etapa ainda podia ser
inspecionada, uma cadeia do tipo:

\begin{lstlisting}
simbolo
-> token
-> representacao numerica
-> transformacao
-> resultado
-> erro
-> atribuicao do erro
-> atualizacao
-> comportamento posterior
\end{lstlisting}

O avanço conceitual relevante para TinyKernel não é a presença de Transformer,
atenção, embedding, softmax ou SGD. Todos esses elementos podem ser realizações
contingentes.

A questão mais profunda é:

\begin{center}
\textbf{qual é o fechamento causal mínimo capaz de converter experiência em mudança
persistente de comportamento?}
\end{center}

Uma abstração candidata é:

\[
\text{estado}_t
\rightarrow
\text{ação}
\rightarrow
\text{consequência observável}
\rightarrow
\text{diferença}
\rightarrow
\text{atualização}
\rightarrow
\text{estado}_{t+1}.
\]

Isso não define aprendizado. Define apenas uma estrutura candidata a ser testada.

\subsection{TinyLogicVision: o mundo entra no circuito}

TinyLogicVision introduziu uma dificuldade adicional: antes de perguntar o que o
modelo aprendeu, tornou-se necessário perguntar \emph{o que exatamente foi
observado e representado}.

A cadeia passou a admitir algo como:

\begin{lstlisting}
fenomeno
-> observacao
-> amostragem
-> suporte
-> representacao
-> entrada
-> transformacao
-> decisao
-> projecao
\end{lstlisting}

Dessa prática emergiram distinções importantes:

\begin{lstlisting}
SUPPORT != DECISION POINT
DECISION POINT != DISPLAY CELL
DISPLAY CELL != H3 CELL
RGB PREVIEW != MULTISPECTRAL TENSOR
H3 != RASTER
\end{lstlisting}

Essas desigualdades sugerem um princípio central para TinyKernel:

\begin{thesisbox}[Princípio de não-colapso]
Duas coisas computacionalmente intercambiáveis em uma implementação não são
necessariamente semanticamente equivalentes no fenômeno investigado.

Um sistema pode continuar executando e, ainda assim, ter perdido a causalidade ou o
significado que legitimava sua saída.
\end{thesisbox}

% =============================================================================
\section{Hipótese central: minimalidade causal}

Considere:

\begin{itemize}
  \item \(P\): um fenômeno explicitamente definido;
  \item \(C\): o contexto no qual o fenômeno é observado;
  \item \(R\): uma realização candidata;
  \item \(W\): um conjunto de testemunhos (\emph{witnesses});
  \item \(\mathcal{I}\): uma família de intervenções permitidas.
\end{itemize}

A notação inicial \(K(P,C,W)\) é útil como intuição, mas deve ser refinada:
o que observamos diretamente não é o fenômeno em si; observamos respostas dos
testemunhos sobre uma realização submetida a um contexto e a intervenções.

Definimos provisoriamente:

\[
W(R,C) = \text{evidência produzida pelos testemunhos sobre }R\text{ em }C.
\]

Uma realização \(R\) é \textbf{suficiente sob o protocolo} quando os testemunhos
previamente definidos continuam sustentando a presença de \(P\) no contexto \(C\).

\[
\operatorname{Suf}(R;P,C,W) = 1.
\]

Para um elemento \(e\) removível de \(R\), definimos uma intervenção subtrativa:

\[
\Delta^{-}_{e}(R) = R \setminus e.
\]

Uma primeira noção de necessidade observada seria:

\[
\operatorname{Nec}(e \mid R;P,C,W)
=
1
\quad\text{se}\quad
\operatorname{Suf}(R)=1
\;\land\;
\operatorname{Suf}(\Delta^{-}_{e}(R))=0.
\]

\subsection{Uma correção importante: irredutível não é necessariamente mínimo}

A expressão

\[
\forall e\in R,\qquad
\operatorname{Suf}(R\setminus e)=0
\]

mostra que \(R\) é \textbf{1-irredutível}: nenhuma remoção unitária preserva o
fenômeno sob o protocolo.

Ela \textbf{não prova} que \(R\) seja a menor realização possível.

Pode existir outra realização \(R'\) com menos elementos, ou uma realização com
elementos diferentes, que preserve o mesmo fenômeno.

Portanto TinyKernel deve distinguir pelo menos:

\begin{description}
  \item[Irredutibilidade local:] nenhuma intervenção elementar autorizada preserva o
  fenômeno.
  \item[Minimalidade relativa:] não encontramos realização menor dentro do espaço de
  busca examinado.
  \item[Minimalidade global:] nenhuma realização possível é menor.
\end{description}

A terceira é, em geral, uma alegação muito mais forte e talvez inalcançável.

\begin{falsifybox}[Regra epistemológica]
TinyKernel não deve chamar uma realização de ``Kernel mínimo'' quando a evidência
apenas demonstra irredutibilidade local.
\end{falsifybox}

% =============================================================================
\section{O problema do witness: não observar não significa que deixou de existir}

Esta é uma das armadilhas mais importantes do programa.

Se uma intervenção faz o witness deixar de detectar o fenômeno, existem pelo menos
duas interpretações:

\begin{enumerate}
  \item o fenômeno realmente foi destruído;
  \item a intervenção destruiu a capacidade do witness de observá-lo.
\end{enumerate}

Logo:

\[
W(P)=0 \;\not\Rightarrow\; P=0.
\]

Da mesma forma:

\[
W(P)=1 \;\not\Rightarrow\; P=1
\]

sem que tenhamos alguma confiança na validade do witness.

Isso significa que o witness não é um acessório do experimento. Ele é parte do
aparato epistemológico e também deve estar sujeito a teste.

\subsection{Triangulação mínima}

Quando possível, o laboratório deverá empregar testemunhos de naturezas distintas:

\begin{itemize}
  \item \textbf{witness operacional}: o sistema produz ou não uma saída?
  \item \textbf{witness causal}: a saída ainda depende da cadeia causal declarada?
  \item \textbf{witness semântico}: a saída ainda significa aquilo que o experimento
  afirma que significa?
  \item \textbf{witness temporal}: a propriedade sobrevive a uma transição de estado
  ou a uma interação posterior?
\end{itemize}

A existência de vários witnesses não produz verdade automática. Ela reduz a chance
de confundir uma falha de observação com uma falha do fenômeno.

% =============================================================================
\section{Quatro tipos de ruptura}

A versão inicial da ideia distinguia três rupturas: operacional, causal e semântica.
A análise do papel do witness exige uma quarta.

\begin{longtable}{>{\bfseries}p{3.4cm}p{10.4cm}}
\toprule
Ruptura & Definição operacional candidata \\
\midrule
\endhead
Operacional &
A realização deixa de produzir o comportamento ou artefato observável requerido
pelo fenômeno. \\

Causal &
A saída continua existindo, mas já não decorre da cadeia causal que fundamentava a
alegação. O sistema ``funciona'', porém por outra razão. \\

Semântica &
A saída computacional permanece, mas uma distinção necessária para seu significado
foi perdida. O valor existe; a interpretação que o legitimava, não. \\

Observacional / evidencial &
O fenômeno pode permanecer, mas a intervenção compromete o mecanismo utilizado para
detectá-lo, medi-lo ou distingui-lo. A ausência observada torna-se inconclusiva. \\
\bottomrule
\end{longtable}

\begin{thesisbox}[Consequência]
\textbf{PASS de software não é witness suficiente de preservação de Kernel.}
Uma intervenção pode preservar execução e destruir causalidade; preservar causalidade
e destruir semântica; ou preservar o fenômeno e destruir somente nossa capacidade
de observá-lo.
\end{thesisbox}

% =============================================================================
\section{O operador experimental fundamental: subtração}

TinyLogicLM e TinyLogicVision foram construídos predominantemente por
\textbf{materialização incremental}: adicionávamos somente aquilo que uma pergunta
exigia.

TinyKernel acrescenta deliberadamente o movimento inverso:

\[
\boxed{
\text{construir}
\;\longleftrightarrow\;
\text{remover}
}
\]

A unidade experimental básica é:

\begin{center}
\begin{tikzpicture}[
  node distance=8mm,
  box/.style={draw, rounded corners, align=center, minimum width=4.6cm,
              minimum height=9mm, fill=TKGray},
  arr/.style={-{Latex[length=2.4mm]}, thick, draw=TKBlue}
]
\node[box, fill=TKBlue!8] (p) {\textbf{Fenômeno}\\definição prévia};
\node[box, below=of p] (r) {\textbf{Realização candidata}};
\node[box, below=of r] (w0) {\textbf{Witness basal}\\fenômeno sustentado?};
\node[box, below=of w0, fill=TKGold!9] (i) {\textbf{Intervenção subtrativa}\\remover / colapsar / substituir};
\node[box, below=of i] (w1) {\textbf{Witness pós-intervenção}};
\node[box, below left=10mm and 3mm of w1, fill=TKGreen!9] (keep) {preserva\\\(\Rightarrow\) dispensável candidato};
\node[box, below right=10mm and 3mm of w1, fill=TKRed!8] (break) {rompe\\\(\Rightarrow\) necessário candidato};

\draw[arr] (p) -- (r);
\draw[arr] (r) -- (w0);
\draw[arr] (w0) -- (i);
\draw[arr] (i) -- (w1);
\draw[arr] (w1) -- (keep);
\draw[arr] (w1) -- (break);
\end{tikzpicture}
\end{center}

A palavra \textbf{candidato} é obrigatória em ambos os ramos.

Remover algo sem observar perda não prova inutilidade universal.
Romper algo após uma remoção não prova necessidade universal.

Cada resultado permanece limitado por \(P\), \(C\), \(W\) e \(\mathcal{I}\).

% =============================================================================
\section{Interações: por que remover uma coisa de cada vez pode enganar}

Considere uma realização com dois mecanismos redundantes, \(a\) e \(b\).

\[
R=\{a,b\}
\]

Se retirarmos \(a\), \(b\) preserva o fenômeno.
Se retirarmos \(b\), \(a\) preserva o fenômeno.

Testes unitários de ablação poderiam sugerir:

\[
a\text{ é dispensável}
\qquad\text{e}\qquad
b\text{ é dispensável}.
\]

Mas:

\[
R\setminus\{a,b\}
\]

pode destruir completamente o fenômeno.

Portanto necessidade pode ser \textbf{conjunta}, e redundância pode esconder
causalidade.

Para uma intervenção sobre um subconjunto \(S\) de elementos:

\[
\Delta^{-}_{S}(R)=R\setminus S.
\]

TinyKernel deverá começar com ablações simples, mas preservar desde o início a
possibilidade de:

\begin{itemize}
  \item remoções pareadas;
  \item remoções de grupos;
  \item substituições por realizações alternativas;
  \item colapso deliberado de distinções;
  \item mudança da ordem das transformações;
  \item intervenções temporais.
\end{itemize}

\begin{markerbox}[TK-MARKER-001 — Minimalidade combinatória | OPEN]
Uma realização pode parecer irredutível ou dispensável dependendo da granularidade
das intervenções. Investigar como distinguir dependência elementar, dependência
conjunta e redundância causal sem transformar o laboratório prematuramente em um
buscador combinatório geral.
\end{markerbox}

% =============================================================================
\section{Kernel não é implementação}

Considere:

\begin{lstlisting}
Realizacao A:
a -> b -> c

Realizacao B:
x -> y
\end{lstlisting}

Se ambas preservam o mesmo fenômeno sob o mesmo contexto e sobrevivem aos mesmos
testes causais relevantes, talvez estejamos observando duas realizações de uma
estrutura mais abstrata.

Isso sugere uma relação de equivalência experimental candidata:

\[
R_1 \sim_{P,C,W,\mathcal{I}} R_2
\]

quando \(R_1\) e \(R_2\) são indistinguíveis \emph{para o fenômeno e protocolo
explicitados}, não em sentido absoluto.

Essa equivalência é deliberadamente local.

Duas realizações podem ser equivalentes para uma pergunta e profundamente diferentes
para outra.

\subsection{O espaço de realizações minimais}

Pode haver:

\[
K_1=\{A,B,C\}
\qquad\text{e}\qquad
K_2=\{A,X\},
\]

ambos suficientes e localmente irredutíveis.

Nesse caso, o objeto de interesse deixa de ser ``o Kernel'' e passa a ser:

\begin{center}
\textbf{o espaço de realizações causalmente minimais de um fenômeno sob um protocolo.}
\end{center}

Esse espaço pode conter famílias, equivalências, redundâncias, compensações e
múltiplas formas de realizar a mesma capacidade.

\begin{markerbox}[TK-MARKER-002 — Pluralidade de Kernel | OPEN]
Testar se fenômenos simples admitem mais de uma realização irredutível não isomórfica.
Se admitirem, abandonar qualquer pressuposto de unicidade do Kernel.
\end{markerbox}

% =============================================================================
\section{Minimalidade causal versus minimalidade descritiva}

Um sistema pode ter uma descrição curta e ainda conter causalidade supérflua.
Outro pode exigir mais símbolos para descrever a mesma estrutura causal.

TinyKernel não deve confundir:

\begin{itemize}
  \item menor número de linhas de código;
  \item menor número de classes;
  \item menor número de bytes;
  \item menor número de estados;
  \item menor descrição textual;
  \item menor estrutura causal suficiente.
\end{itemize}

Essas medidas podem ser correlacionadas em alguns experimentos, mas não são
sinônimos.

\begin{markerbox}[TK-MARKER-003 — Unidade de complexidade | OPEN]
Que medida, se alguma, permite comparar a complexidade de realizações diferentes sem
reificar escolhas de implementação? Contagem de componentes é suficiente apenas para
experimentos extremamente controlados.
\end{markerbox}

% =============================================================================
\section{O primeiro fenômeno não deve se chamar “aprendizado”}

A formulação inicial de TK-0001 usava ``persistência de aprendizado''. Isso carrega
um problema: a palavra \emph{aprendizado} já embute uma interpretação.

Para reduzir circularidade, o primeiro fenômeno deve ser mais fraco e diretamente
observável.

\begin{thesisbox}[Fenômeno TK-0001 — formulação recomendada]
\textbf{Persistência adaptativa:}
um sistema sofre uma experiência \(E\), altera algum estado interno em função dessa
experiência e, em uma interação posterior na qual \(E\) não é reapresentada, produz
um comportamento diferente do baseline em consequência daquela alteração persistida.
\end{thesisbox}

Isso não prova que o sistema ``aprendeu''.

Permite perguntar depois:

\begin{center}
\emph{Que condições adicionais seriam necessárias para distinguir persistência
adaptativa de memorização, ajuste, condicionamento ou aprendizado?}
\end{center}

Essa é precisamente a lógica do laboratório: começar com um fenômeno pequeno e
deixar conceitos mais fortes emergirem apenas se a diferença se tornar necessária.

% =============================================================================
\section{TK-0000 e TK-0001: materialidade inicial}

\subsection{TK-0000 — Bootstrap do aparato}

TK-0000 não deve buscar um Kernel substantivo.

Seu único objetivo é provar que o laboratório consegue executar de forma
reproduzível:

\begin{lstlisting}
baseline
-> witness
-> intervention
-> witness
-> classification
\end{lstlisting}

Critério de saída:

\begin{itemize}
  \item build limpo;
  \item execução determinística;
  \item witness basal explícito;
  \item uma intervenção subtrativa trivial;
  \item restauração do baseline;
  \item resultado reproduzível;
  \item nenhum claim sobre minimalidade causal.
\end{itemize}

\subsection{TK-0001 — Persistência adaptativa}

Uma realização inicial pode ser propositalmente banal:

\[
s_0=0.
\]

O sistema recebe uma experiência com um alvo \(y\in\{0,1\}\), produz uma ação baseada
em \(s\), recebe informação sobre a diferença e atualiza \(s\). Em uma interação
posterior, sem reapresentar a experiência original, o comportamento deve refletir o
estado alterado.

A realização candidata pode conter:

\begin{lstlisting}
estado
-> entrada
-> acao
-> feedback
-> diferenca
-> atualizacao
-> estado persistido
-> interacao posterior
\end{lstlisting}

Intervenções iniciais candidatas:

\begin{longtable}{p{4.1cm}p{5.0cm}p{5.0cm}}
\toprule
Intervenção & Previsão antes da execução & O que a observação pode distinguir \\
\midrule
\endhead
Remover atualização &
Experiência ocorre, mas não altera estado &
Se mudança persistente exige atualização nesta realização \\

Remover persistência &
Mudança pode ocorrer durante a interação, mas desaparece depois &
Se o fenômeno definido exige continuidade temporal \\

Remover feedback &
Mudança pode ainda ocorrer por outra regra &
Se feedback é necessário para persistência adaptativa ou apenas para uma classe
específica de adaptação \\

Substituir feedback por ruído &
O estado pode mudar sem se orientar pela experiência &
Se orientação por consequência precisa ser distinguida de mera mutação \\

Colapsar estado e ação &
Pode preservar saída em um instante e perder memória causal &
Se estado independente é necessário para explicar persistência \\
\bottomrule
\end{longtable}

\begin{falsifybox}[Proibição de interpretação retroativa]
Se TK-0001 mostrar que feedback não é necessário para o fenômeno definido, o resultado
não deve ser ``corrigido'' redefinindo silenciosamente o fenômeno para fazer feedback
parecer necessário.

A definição prévia permanece registrada; uma definição sucessora deve receber outro
experimento.
\end{falsifybox}

% =============================================================================
\section{Classificação experimental de resultados}

Cada intervenção deverá produzir uma classificação limitada ao protocolo.

\begin{description}
  \item[PRESERVED:] os witnesses preregistrados sustentam que o fenômeno permanece.
  \item[BROKEN\_OPERATIONAL:] o fenômeno falha no nível operacional.
  \item[BROKEN\_CAUSAL:] a saída permanece, mas sua dependência causal requerida foi
  quebrada.
  \item[BROKEN\_SEMANTIC:] a saída permanece, mas a distinção semântica necessária
  foi perdida.
  \item[WITNESS\_COMPROMISED:] a intervenção impede conclusão porque afeta o aparato
  de observação.
  \item[INCONCLUSIVE:] os dados não distinguem as alternativas definidas.
\end{description}

Nenhuma dessas classificações equivale automaticamente a:

\begin{center}
\texttt{KERNEL\_FOUND}
\end{center}

Essa etiqueta não deve existir nos primeiros experimentos.

% =============================================================================
\section{Da visão ao TinyKernel: intervenções semânticas futuras}

TinyLogicVision sugere uma família especialmente fértil de experimentos:
\textbf{colapso de distinções}.

Em vez de apenas apagar componentes, podemos tornar duas categorias
indistinguíveis e observar se a alegação científica continua legítima.

Exemplos candidatos:

\begin{lstlisting}
support := decision_point
preview_rgb := multispectral_tensor
h3_cell := raster_support
display_cell := observation_support
\end{lstlisting}

O objetivo não será verificar se o software continua executando.

Será verificar se:

\begin{enumerate}
  \item a cadeia causal continua reconstruível;
  \item a saída continua semanticamente atribuível à observação declarada;
  \item dois mundos materialmente diferentes continuam distinguíveis.
\end{enumerate}

\begin{markerbox}[TK-MARKER-004 — Colapso semântico | OPEN]
Formalizar uma classe de intervenção em que nenhuma função é removida, mas duas
distinções são fundidas. Testar se a perda de distinguibilidade pode servir como
evidência de necessidade ontológica relativa ao domínio.
\end{markerbox}

% =============================================================================
\section{Da minimalidade ao Sister-Kernel}

TinyKernel não deve nascer como laboratório de validação do Sister-Kernel.

O fluxo correto é unidirecional apenas quando existe evidência suficiente:

\begin{center}
\begin{tikzpicture}[
  node distance=10mm and 17mm,
  box/.style={draw, rounded corners, align=center, minimum width=3.8cm,
              minimum height=10mm, fill=TKGray},
  arr/.style={-{Latex[length=2.5mm]}, thick, draw=TKBlue}
]
\node[box, fill=TKBlue!8] (lm) {TinyLogicLM\\capacidade aprendida};
\node[box, right=of lm, fill=TKTeal!8] (vis) {TinyLogicVision\\observação e inferência};
\node[box, below=14mm of $(lm)!0.5!(vis)$, fill=TKGold!10] (tk) {\textbf{TinyKernel}\\minimalidade causal};
\node[box, below=of tk] (ev) {evidência experimental\\limitada ao protocolo};
\node[box, below=of ev, fill=TKGreen!8] (sis) {Sister-Kernel / Atmos /\\Nexo / Praxis / outros};

\draw[arr] (lm) -- (tk);
\draw[arr] (vis) -- (tk);
\draw[arr] (tk) -- (ev);
\draw[arr] (ev) -- node[right,align=left]{\small somente após\\\small decisão própria} (sis);
\end{tikzpicture}
\end{center}

Quando o laboratório estiver maduro o suficiente, poderá testar materialmente
distinções como:

\begin{lstlisting}
Participant != Role
Role != Authority
Evidence != Decision
Binding != Authority
Artifact != Provenance
\end{lstlisting}

A estratégia será criar situações que \emph{devem permanecer distinguíveis} e testar
se uma representação que colapsa a distinção perde capacidade de responder a uma
pergunta operacional relevante.

Se perder, isso será evidência em favor de uma distinção necessária \emph{naquele
escopo}. Não será uma licença automática para universalizar a entidade.

% =============================================================================
\section{“Sempre pronto e sempre incompleto” como propriedade verificável}

A frase deixa de ser apenas filosofia de desenvolvimento e passa a ser requisito
operacional.

\subsection{Sempre pronto}

Um estado publicável do repositório deve satisfazer, no mínimo:

\begin{lstlisting}
HEAD compila
HEAD executa
HEAD verifica
experimentos publicados sao reproduziveis
estado atual possui witness
claims atuais sao explicitos
limites de evidencia sao explicitos
\end{lstlisting}

\subsection{Sempre incompleto}

O mesmo estado deve preservar:

\begin{lstlisting}
nenhuma ontologia final declarada
nenhum Kernel universalizado alem da evidencia
novos contextos podem refutar minimalidade anterior
realizacoes alternativas permanecem possiveis
resultados inconclusivos permanecem resultados
a proxima pergunta permanece aberta
\end{lstlisting}

Logo:

\[
\boxed{\text{READY}\neq\text{COMPLETE}}
\]

\begin{thesisbox}[Invariante de projeto]
Um release pode ser tecnicamente fechado e epistemicamente aberto.

Completude operacional não autoriza completude ontológica.
\end{thesisbox}

% =============================================================================
\section{Constituição mínima de TinyKernel}

Esta constituição deve ser pequena o suficiente para ser preservada e forte o
suficiente para impedir que o laboratório se transforme em uma máquina de confirmar
suas próprias suposições.

\begin{enumerate}[label=\textbf{P\arabic*.}]
  \item \textbf{Fenômeno antes de componentes.}
  Nenhum componente é considerado necessário antes de o fenômeno investigado ser
  definido de forma observável.

  \item \textbf{Nenhuma necessidade por tradição.}
  Práticas comuns, padrões arquiteturais e popularidade de uma tecnologia não
  constituem evidência de pertencimento a Kernel.

  \item \textbf{Necessidade exige intervenção.}
  Alegações de necessidade devem, sempre que materialmente possível, sobreviver a
  testes subtrativos, substitutivos ou de colapso.

  \item \textbf{Witness não é fenômeno.}
  Falha de observação não deve ser confundida automaticamente com ausência do
  fenômeno.

  \item \textbf{Execução não basta.}
  Funcionamento computacional não é evidência suficiente se causalidade ou semântica
  forem destruídas.

  \item \textbf{Irredutível não significa globalmente mínimo.}
  O laboratório deve declarar a força exata da minimalidade demonstrada.

  \item \textbf{Implementação não é Kernel.}
  Realizações distintas podem satisfazer a mesma estrutura causal relevante.

  \item \textbf{Unicidade não é presumida.}
  Pode existir mais de uma realização minimal suficiente para o mesmo fenômeno.

  \item \textbf{Necessidade pode ser conjunta.}
  Interações e redundâncias podem tornar ablações unitárias insuficientes.

  \item \textbf{Resultados são contextuais.}
  Toda alegação é limitada pelo fenômeno, contexto, witnesses e família de
  intervenções examinados.

  \item \textbf{História experimental é preservada.}
  Previsões, falhas, resultados negativos e inconclusões não podem ser reescritos
  retroativamente para acomodar estados sucessores.

  \item \textbf{O conceito de Kernel é falsificável.}
  Se ``Kernel'' deixar de ser uma categoria útil diante da evidência, o laboratório
  deve continuar existindo sem ela.
\end{enumerate}

% =============================================================================
\section{Estrutura inicial do repositório}

A estrutura deve permanecer deliberadamente austera:

\begin{lstlisting}
TinyKernel/
|-- CMakeLists.txt
|-- README.md
|-- bin/
|   `-- tinykernel
|-- src/
|-- tests/
`-- experiments/
    |-- TK-0000/
    |   |-- question.md
    |   |-- preregistration.md
    |   `-- result.md
    `-- TK-0001/
        |-- question.md
        |-- preregistration.md
        `-- result.md
\end{lstlisting}

Inicialmente, não introduzir:

\begin{itemize}
  \item banco de dados;
  \item servidor web;
  \item sistema de plugins;
  \item framework de experimentos genérico;
  \item DSL de Kernel;
  \item grafo universal;
  \item ontologia formal;
  \item dependência de LLM;
  \item integração direta com Sister-Kernel.
\end{itemize}

O primeiro executável pode precisar somente de:

\begin{lstlisting}
./bin/tinykernel verify
./bin/tinykernel experiment TK-0000
./bin/tinykernel experiment TK-0001
\end{lstlisting}

Qualquer nova superfície deve responder a uma necessidade observada.

% =============================================================================
\section{Programa experimental inicial}

\begin{longtable}{p{2.3cm}p{4.6cm}p{7.3cm}}
\toprule
Experimento & Objeto & Pergunta \\
\midrule
\endhead
TK-0000 &
Bootstrap do aparato &
Conseguimos executar baseline, intervenção e witness de forma determinística sem
introduzir claims substantivos? \\

TK-0001 &
Persistência adaptativa &
Quais relações são necessárias para que uma experiência produza mudança persistente
de comportamento em uma interação posterior? \\

TK-0002 &
Witness e observabilidade &
Podemos destruir o witness sem destruir o fenômeno? Como distinguir ausência de
evidência de evidência de ausência? \\

TK-0003 &
Realizações alternativas &
Duas máquinas estruturalmente diferentes podem preservar o mesmo fenômeno e resistir
a intervenções causalmente equivalentes? \\

TK-0004 &
Redundância causal &
Quais necessidades só aparecem em remoções conjuntas e permanecem invisíveis em
ablações unitárias? \\

TK-0010 &
Colapso semântico &
Uma implementação pode continuar funcional após fundir duas distinções e, ainda
assim, perder o significado necessário para sua alegação? \\

TK-0100 &
Primitivas federadas &
Participant, Role, Authority, Evidence e Decision são distinções experimentalmente
necessárias em casos controlados de composição? \\
\bottomrule
\end{longtable}

A numeração deliberadamente deixa espaço entre os primeiros experimentos e os
experimentos de sistemas federados. A distância é metodológica: TinyKernel deve
provar seu aparato em fenômenos pequenos antes de tentar explicar o SisTer.

% =============================================================================
\section{Marcadores de investigação}

Os marcadores abaixo são perguntas abertas, não backlog de implementação.

\begin{markerbox}[TK-MARKER-005 — Fronteira do fenômeno | OPEN]
Quando duas definições muito próximas devem ser tratadas como o mesmo fenômeno em
contextos diferentes, e quando a mudança de definição exige um novo experimento?
\end{markerbox}

\begin{markerbox}[TK-MARKER-006 — Causalidade sem intervenção total | OPEN]
Há situações em que uma relação causal relevante não pode ser removida de modo
seguro ou semanticamente limpo. Que evidência substitutiva seria aceitável sem
rebaixar ``necessidade'' a mera plausibilidade?
\end{markerbox}

\begin{markerbox}[TK-MARKER-007 — Granularidade | OPEN]
Um ``elemento'' removível pode ser função, estado, relação, ordem, fronteira,
informação ou distinção semântica. Como evitar que a granularidade escolhida determine
artificialmente o Kernel encontrado?
\end{markerbox}

\begin{markerbox}[TK-MARKER-008 — Temporalidade | OPEN]
Alguns fenômenos só existem como trajetória. Um Kernel temporal pode depender não
apenas de estados, mas da ordem legítima das transições? Investigar kernels de
trajetória, não apenas de configuração.
\end{markerbox}

\begin{markerbox}[TK-MARKER-009 — Contexto e transferência | OPEN]
Que parte de uma estrutura minimal descoberta em \(C_1\) sobrevive em \(C_2\)?
Definir transferência como nova hipótese, nunca como consequência automática.
\end{markerbox}

\begin{markerbox}[TK-MARKER-010 — O fim do Kernel | OPEN]
Que resultado experimental faria abandonar a palavra ``Kernel''? Registrar desde o
início condições sob as quais a categoria se mostraria enganosa, excessivamente
contextual ou não identificável.
\end{markerbox}

% =============================================================================
\section{Critérios para transferência a outros projetos}

Nenhum resultado de TinyKernel deve ser transferido para Sister-Kernel, Nexo, Atmos,
Praxis, Memória, Infra ou outro projeto apenas por analogia.

Uma transferência deve declarar:

\begin{enumerate}
  \item fenômeno original estudado;
  \item contexto original;
  \item witnesses utilizados;
  \item intervenções realizadas;
  \item força real da alegação obtida;
  \item novo fenômeno e novo contexto de destino;
  \item quais condições de equivalência são assumidas;
  \item qual experimento de destino pode falsificar a transferência.
\end{enumerate}

\begin{warningbox}[Regra de transferência]
\textbf{Analogia gera hipótese; não gera autoridade.}
\end{warningbox}

% =============================================================================
\section{Definições operacionais candidatas — v0.1.0}

Estas definições existem apenas para tornar os primeiros experimentos formuláveis.

\begin{description}
  \item[Fenômeno:] propriedade, capacidade, relação ou trajetória que o experimento
  pretende distinguir como presente ou ausente.

  \item[Contexto:] condições materiais, temporais e observacionais dentro das quais a
  alegação sobre o fenômeno é válida.

  \item[Realização:] configuração concreta capaz de produzir observações relacionadas
  ao fenômeno.

  \item[Witness:] mecanismo explicitado que produz evidência sobre uma propriedade da
  realização; não é sinônimo do fenômeno.

  \item[Intervenção:] alteração controlada feita para distinguir explicações ou testar
  necessidade, suficiência, causalidade ou semântica.

  \item[Subtração:] intervenção que remove uma parte, relação, estado, etapa ou
  distinção de uma realização.

  \item[Colapso:] intervenção que funde categorias antes distinguíveis sem
  necessariamente remover computação.

  \item[Suficiência relativa:] preservação do fenômeno segundo witnesses e contexto
  preregistrados.

  \item[Necessidade relativa:] perda do fenômeno após intervenção comparável, dentro
  do protocolo examinado.

  \item[Irredutibilidade local:] inexistência de uma intervenção elementar autorizada
  que preserve o fenômeno.

  \item[Kernel candidato:] realização ou estrutura causal que, após uma sequência de
  experimentos, possui evidência suficiente para ser tratada como mínima em um escopo
  explicitado. O termo não deve aparecer como resultado de TK-0000.
\end{description}

% =============================================================================
\section{O que este documento não demonstra}

Este documento não demonstra:

\begin{itemize}
  \item que todo fenômeno possui um Kernel identificável;
  \item que um Kernel é único;
  \item que um Kernel pode ser encontrado por ablação;
  \item que minimalidade causal é computável em geral;
  \item que witnesses disponíveis conseguem distinguir preservação de ruptura;
  \item que a noção de Kernel transfere entre domínios;
  \item que TinyKernel produzirá resultados úteis ao SisTer;
  \item que persistência adaptativa é aprendizado;
  \item que causalidade pode ser reduzida a uma lista de componentes;
  \item que a constituição P1--P12 permanecerá intacta.
\end{itemize}

O valor deste documento está justamente em tornar essas incertezas explícitas antes
que a implementação as esconda.

% =============================================================================
\section{Primeiro estado autorizado}

Após publicação deste documento, o único próximo estado tecnicamente justificável é:

\begin{center}
\boxed{
\text{documentar TK-0000 antes de implementar TK-0001}
}
\end{center}

TK-0000 deverá provar o aparato, não a teoria.

Somente depois de TK-0000 reproduzível, TK-0001 poderá perguntar sobre persistência
adaptativa.

Nenhuma implementação deve tentar resolver antecipadamente TK-0002, TK-0003 ou os
marcadores abertos.

% =============================================================================
\section{Síntese final}

O TinyKernel nasce de uma inversão metodológica.

Durante muito tempo, construir sistemas significou perguntar:

\begin{center}
\emph{“o que devemos acrescentar para que isto funcione?”}
\end{center}

O novo laboratório acrescenta uma pergunta complementar:

\begin{center}
\Large\bfseries
“o que podemos retirar antes que isto deixe de ser isto?”
\end{center}

Essa pergunta parece simples, mas obriga a distinguir funcionamento, causalidade,
significado, observação, contexto e realização.

Talvez o Kernel seja uma estrutura causal mínima.
Talvez existam muitos Kernels para o mesmo fenômeno.
Talvez minimalidade dependa fortemente do contexto.
Talvez a própria palavra Kernel seja inadequada.

Todas essas possibilidades permanecem abertas.

\begin{thesisbox}[Compromisso fundador]
TinyKernel não será construído para provar que a teoria de TinyKernel está correta.

Será construído para tornar materialmente possível descobrir onde ela está errada.
\end{thesisbox}

\vspace{0.8cm}

\begin{center}
{\sffamily\Large\bfseries\color{TKBlue}
Sempre pronto. Sempre incompleto.}
\end{center}

% =============================================================================
\appendix

\section{Registro de proveniência do documento}

\begin{tabularx}{\textwidth}{>{\bfseries}p{4.1cm}X}
\toprule
Evento de origem & Diálogo exploratório sobre a essência do Kernel e os avanços
conceituais derivados de TinyLogicLM e TinyLogicVision \\
Data da consolidação & 09/09/2026 \\
Transformação realizada & Conversão de diálogo em documento fundador com
formalizações, limites, constituição e programa experimental \\
Elementos preservados & minimalidade causal; construir/remover; \(K(P,C,W)\);
suficiência; irredutibilidade; múltiplas realizações; rupturas operacional, causal e
semântica; ``sempre pronto e sempre incompleto'' \\
Elementos adicionados nesta versão & distinção entre 1-irredutibilidade e minimalidade
global; papel falível do witness; ruptura observacional; necessidade conjunta;
relação de equivalência experimental; TK-0000; reformulação de TK-0001 como
persistência adaptativa; constituição P1--P12; marcadores de investigação \\
Natureza das adições & hipóteses e refinamentos conceituais, não resultados \\
Próxima revisão esperada & após preregistro e execução de TK-0000 \\
\bottomrule
\end{tabularx}

\section{Histórico de versões}

\begin{longtable}{p{2.2cm}p{2.8cm}p{9.4cm}}
\toprule
Versão & Data & Mudança \\
\midrule
\endhead
0.1.0 & 09/09/2026 &
Primeiro registro fundador. Formaliza minimalidade causal como objeto de pesquisa,
separa fenômeno de witness, distingue irredutibilidade local de minimalidade,
introduz ruptura observacional, define constituição P1--P12 e programa experimental
TK-0000--TK-0100. \\
\bottomrule
\end{longtable}

\end{document}
